\scanpage{9}
[ \uline{NB}. $\prod_{k'/k}G_{k'}$ est de t. m. ] donc grâce à $H^{n+2}=0$ sur les groupes de t. m. on voit que
\diagram{\begin{tikzcd}[ampersand replacement=\&,row sep=0.25em]
H^{n+1}(\prod_{k'/k}G_{k'})\arrow[r] \& H^{n+1}(G)\quad\text{est surjectif,}\\
\rotatebox{90}{$\simeq$}\\
H^{n+1}(G_{k'})
\end{tikzcd}}
donc pour vérifier que $H^{n+1}(G)=0$, il suffit de le faire pour $H^{n+1}(G_{k'})$. Cela nous ramène au cas où $G$ est ``\unsure{splitté|splité}'' par une $p$-extension, \uline{ce qui implique que} $G$ \uline{se dévisse en groupes isomorphes à} $\mu_p$. Donc on est ramené à vérifier que $H^{n+1}(\mu_p)=0$, ce qui se voit sur la suite exacte
\[
0\longrightarrow\mu_p\longrightarrow G_m\longrightarrow G_m\longrightarrow 0
\]

\uline{Lemme} 3 \deleted{\unsure{S}} On a
\[
\operatorname{cd}^{\substack{\rm \unsure{plat|rad}\\\rm add}}_p(k)=
\begin{cases}0&\text{si }k\text{ parfait}\\1&\text{si }k\text{ imparfait}\end{cases}
\]

En effet, un groupe radiciel additif se dévisse en groupes $\alpha_p$ \unsure{\ldots\,.}
% archive sheet number: 8

